\documentclass[10pt,a4paper]{amsart}
\usepackage[margin=19mm]{geometry}
\usepackage{amsmath,amssymb,mathtools}
\usepackage[scr=rsfs,cal=euler]{mathalfa}
\usepackage{xfrac}
\usepackage[unicode,hidelinks,bookmarks=false]{hyperref}
\newcommand{\ie}{i.e.\ }
\newcommand{\cf}{cf.\ }
\newcommand{\resp}{respectively\ }
\newcommand{\Subsection}{\subsection}
\providecommand{\nobreakdash}{\nobreak\hbox{-}}
\setlength{\emergencystretch}{2em}
\multlinegap0pt
\usepackage{bm}
\usepackage{bbm}
\usepackage{dsfont}
\usepackage{xspace}
\usepackage{cancel}
\usepackage{mathtools}
\usepackage{amsthm}
\makeatletter
\@ifundefined{thm}{\newtheorem{thm}{Thm}}{}
\@ifundefined{lem}{\newtheorem{lem}[thm]{Lemma}}{}
\makeatother
\makeatletter
\newcommand{\inc}{\mathrm{inc}}
\expandafter\ifx\csname proofidea\endcsname\relax
\newenvironment{proofidea}{\noindent \paragraph{\emph{Proof idea.}}}{\hfill $\Box$}
\fi
\expandafter\ifx\csname proofomitted\endcsname\relax
\newenvironment{proofomitted}{\noindent \paragraph{\emph{Proof omitted.}}}{\hfill $\Box$} 
\fi
\newcommand{\ssm}{\smallsetminus}
\makeatother
\title[Hook immanantal inequalities for totally nonnegative matrice]{Hook immanantal inequalities for totally nonnegative matrices}
\author{Mark Skandera}
\date{Source: arXiv:2510.00327v1 (CC BY 4.0)}
\begin{document}
\maketitle

\noindent\emph{Source and licence.} Hook immanantal inequalities for totally nonnegative matrices. Version arXiv:2510.00327v1 (CC BY 4.0), distributed under the Creative Commons Attribution 4.0 International licence, as recorded in the arXiv licence metadata for this exact version and shown on the abstract page https://arxiv.org/abs/2510.00327v1. Full rights notice: licenses/arxiv-2510.00327-CC-BY-4.0.txt.\par\smallskip

\noindent\textit{Editorial scope.} Counted unit: the Lem whose source label is \texttt{l:hookineq} in the article (source0.tex lines 1885--1904, source label \texttt{l:hookineq}), delivered together with the article's own complete original proof of it (source0.tex lines 1905--1973, one \texttt{proof} environment of 69 lines). No mathematics is written, repaired or re-derived: statement and proof are the article's own text line for line. Required context retained uncounted: none. Every definition, display and label of those lines is retained unchanged. Omitted from this package and not redistributed: 1--1884, 1974--3130. Nothing inside a retained excerpt is omitted or altered: every retained body is delivered exactly as the article's lines are written, so no omission receipt of any kind accompanies this package. Citations: the retained excerpts carry no citation command at all, so no bibliography entry of the article is transcribed and no reference is redistributed. Reference adaptation: none was needed and none was made; every reference inside the delivered unit resolves inside it, so no unresolved reference or citation is produced. Printed slips kept literal: no printed word, symbol, hypothesis or number of the counted statement or of its proof was altered. Layout and class: the article's own class is \texttt{amsart} with packages including amssymb, amsmath, amsthm, graphicx, enumerate, tikz, tabularx, tabmac\_avi; the wrapper is the lane's pinned \texttt{amsart} editorial wrapper at 10pt on A4, which restates the theorem-like environments the retained text needs and adds the article's own macro definitions that the retained text needs (\texttt{\textbackslash inc}, \texttt{\textbackslash ssm}), with extra wrapper packages (bm, bbm, dsfont, xspace, cancel, mathtools). The delivered numbering is the unit's own. Assets: no retained span contains a figure environment, a graphics-inclusion command, a PDF-page inclusion, a table or a TikZ picture; no image file is copied into this package, the package declares no asset dependency and no image file is redistributed with it. Licence: version arXiv:2510.00327v1 of this article is distributed under the Creative Commons Attribution 4.0 International licence, as recorded in the arXiv licence metadata for this exact version and shown on the abstract page https://arxiv.org/abs/2510.00327v1. A full rights notice is delivered at licenses/arxiv-2510.00327-CC-BY-4.0.txt. Characters: every retained excerpt is pure ASCII, so the delivered engine, XeLaTeX with its default Latin Modern text font, reports no missing character.
\begin{lem}\label{l:hookineq}
  For each naturally labeled $n$-element poset $P$ and $k = 2,\dotsc,n$
  we have
  %and define $d_i$ to be the number of
  %standard $P$-tableaux of shape $i1^{n-i}$.
  %Then we have $(k-1)d_k \geq (n-k+1)d_{k-1}$.
  \begin{equation}\label{eq:binomapprox}
    (k-1) \chi^{k1^{n-k}}(\inc(P)) \geq (n-k+1) \chi^{k-1,1^{n-k+1}}(\inc(P)).
  \end{equation}
  Furthermore, the difference
%  $(k-1)d_k - (n-k+1)d_{k-1}$
   $(k-1) \chi^{k1^{n-k}}(\inc(P)) - (n-k+1) \chi^{k-1,1^{n-k+1}}(\inc(P))$
  equals the number of standard $P$-tableaux
  of shape $k1^{n-k}$ in which one entry of the first row is marked
  and is incomparable in $P$ to an entry in an earlier column of the tableau.
  %satisfies at least one of the conditions
  %\begin{enumerate}
  %\item $j$ for all $i$ to the left of $j$ in row $1$,
  %  \item $i$ is 
\end{lem}

\begin{proof}
  Let $\mathcal C_k = \mathcal C_k(P)$
  be the set of standard $P$-tableaux of shape
  $k 1^{n-k}$ in which one entry of column $1$
  other than that in position $(1,1)$ is marked.
  Let $\mathcal R_k = \mathcal R_k(P)$
  be the set of standard $P$-tableaux of shape
  $k 1^{n-k}$ in which one entry of row $1$
  other than that in position $(1,1)$ is marked.
  The inequality (\ref{eq:binomapprox}) asserts that
  $|\mathcal R_k| \geq |\mathcal C_{k-1}|$.
  To see this, we define a family of maps
  $\{f_k: \mathcal C_{k-1} \rightarrow \mathcal R_k \,|\, k = 2,\dotsc, n\}$,
  by letting
  $f_k(U)$ be the tableau constructed from $U$ by
  removing the marked element from the first column of $U$,
  and reinserting it as far as possible to the right in row $1$
  subject to the requirement that it be a record.
  %o that it is a record.

  To see that the map $f_k$ is well-defined,
  let $m \in [n]$ be the marked element in column $1$ of $U$.
  Since $U$ is standard,
  $m$ must be greater than the element in position $(1,1)$.
  Thus it will certainly be inserted into positions $(1,2), \dotsc, (1,k-1)$
  or at the end of row $1$, creating a tableau of shape $k1^{n-k}$ with
  one marked element in the first row.
  We claim that the map $f_k$ is injective.
  To see this,
  %that it is injective,
  consider a tableau $V \in \mathcal R_{k}$
  which satisfies $V = f_k(U)$ for some $U \in \mathcal C_{k-1}$.
  A marked element $i$ in row $1$ of $V$ will necessarily be greater in $P$
  than all elements to its left in row $1$, and will be comparable in $P$
  to all elements in column $1$.
  To recover $U$, remove $i$ from row $1$ of $V$
  and insert it into the unique position of column $1$ so that entries there
  increase.  The resulting tableau will still be $P$-semistrict
  because if the element preceding $i$ in $V$, say $h$,
  and the element following $i$ in $V$, say $j$,
  satisfy $h >_P j$, then we have $i >_P h >_P j$, contradicting
  the semistrictness of row $1$ of $V$.
  
  To see that the difference
%  $(k-1)d_k - (n-k+1)d_{k-1}$
  $(k-1) \chi^{k1^{n-k}}(P) - (n-k+1) \chi^{k-1,1^{n-k+1}}(P)$
  has the claimed interpretation, consider the
  elements of $\mathcal R_k \ssm f(C_{k-1})$.
  These are standard $P$-tableaux $V$
  which contain a marked element $i$ in row $1$
  which is not a record in row $1$, or which is incomparable to
  some element of column $1$.  We claim that if $i$ is
  %comparable to all elements of column $1$ and
  not a record in row $1$, then it must
  be incomparable to an element to its left in row $1$.
  Assume otherwise: assume that $i$ follows $h$ in row $1$ with $h <_P i$
  and that some element $j >_P i$ appears earlier than $h$ in row $1$.
  Assume that $j$ is the rightmost such element with these properties.
  Then we have $j >_P h$.
  Let $g$ be the element immediately following $j$ in row $1$.
  By our choice of $j$ we cannot have $g >_P j$,
  and since $V$ is standard we cannot have $j >_P g$.
  Thus $j$ must be incomparable to $g$.
  %we cannot have $j <_P g$ or $g >_P i$.
  %Since $V$ is standard, we cannot have $j >_P g$.
  By our choice of $j$ we cannot have $g >_P i$,
  and our assumption on $i$ we cannot have $g$ incomparable to $i$.
  Thus we have $g <_P i$.  But then we have $g <_P i <_P j$, a contradiction.
\end{proof}

\end{document}
